% !TEX root = ../foundations.tex

Before beginning the technical development, we offer a few examples which illustrate the basic appeal of geometric cochains.
They are intentionally brief; the paper's main work is to justify the constructions that make such examples legitimate.

\subsubsection{Crossing a ray}

Let \(M = \R^2 - \{0\}\), oriented by the ordered basis \((\e_1,\e_2)\).
Let
\[
W = \{t\e_1 \mid t > 0\} \subset M
\]
be the positive \(\e_1\)-axis, co-oriented by the normal vector \(\e_2\).
The inclusion \(W \into M\) is proper, so it defines a geometric cochain even though its domain is noncompact.
Since \(W\) has no boundary, it is a geometric cocycle which represents the usual generator of \(H^1_\Gamma(M) \cong \Z\).
Indeed, let \(S^1 \into M\) be the unit circle with its counterclockwise orientation.
It is transverse to \(W\) and intersects \(W\) once, at \(\e_1\).
The sign of this intersection is positive, since the tangent vector to \(S^1\) at \(\e_1\) is \(\e_2\), which agrees with the chosen normal co-orientation of \(W\).
Thus the pairing of the class of \(W\) with the fundamental class of \(S^1\) is \(1\), and therefore \(W\) represents a generator.

This simple example already displays two useful features of the theory.
First, the representative is integral rather than real-valued.
It records winding number by a discrete crossing count: a positively oriented loop contributes \(+1\) each time it crosses the ray in the positive co-oriented direction.
By contrast, the de Rham representative \((2\pi)^{-1}d\theta\) records the same invariant continuously around the circle and only after extension of coefficients to \(\R\).
One can approximate the crossing representative analytically by normalized angular bump forms supported in narrower and narrower sectors around the positive ray, but the limiting object is naturally a geometric, or current-like, representative.

Second, the domain of a geometric cochain need not be compact.
The correct finiteness condition is properness of the map to the target.
Thus geometric cohomology admits representatives that are natural in pictures but less directly visible in approaches that identify ordinary cohomology with duals of compact cycles.
%Equivalently, this ray is familiar from Borel--Moore or locally finite homology, but the geometric-cochain viewpoint uses it directly as a cohomology representative.

\subsubsection{Cup products}

Once the additive cohomology groups are known, geometric representatives can often be chosen so that cup products are read from intersections.
Projective spaces provide the basic examples.

For complex projective space, let \(h \in H^2_\Gamma(\mathbb{C}P^n)\) be represented by a hyperplane
\[
\mathbb{C}P^{n-1} \into \mathbb{C}P^n,
\]
with its normal co-orientation induced by the complex structure.
More generally, a linear subspace
\[
\mathbb{C}P^{n-k} \into \mathbb{C}P^n
\]
represents \(h^k\).
The product calculation is then the elementary fact that transverse linear subspaces intersect in a linear subspace, with codimensions adding.
Thus the presentation
\[
H^*(\mathbb{C}P^n;\Z) \cong \Z[h]/(h^{n+1})
\]
reflects the fact that \(k\) independent hyperplane conditions cut out \(\mathbb{C}P^{n-k}\), while \(n+1\) independent hyperplane conditions have empty intersection.

\medskip
The real setting illustrates the distinction between orientations and co-orientations, with odd-dimensional projective spaces being orientable but even-codimension projective spaces being co-orientable.

With mod-two coefficients, the linear inclusions
\[
\mathbb{R}P^{n-k} \into \mathbb{R}P^n
\]
represent the powers of the degree-one generator.
With integral coefficients, such an inclusion defines a geometric cochain only when its normal bundle is orientable.
This separates two pieces of structure that coincide too easily in examples over \(\mathbb{F}_2\): an integral chain requires an orientation of its domain, whereas an integral cochain requires a co-orientation of its map to the target.

\medskip
These examples illustrate the basic geometric procedure.
After the additive groups have been computed, one chooses geometric cocycles dual to geometric cycles and then computes products by making the cocycles transverse and intersecting them.

\subsubsection{Characteristic classes}

Characteristic classes also admit geometric representatives of this kind.
Let \(E \to M\) be a rank \(r\) vector bundle.
The zero section
\[
M \into E
\]
is a canonical submanifold of the total space.
With an orientation of \(E\) it represents the Thom class geometrically; without such an orientation it does so with \(\mathbb{F}_2\)-coefficients.
Pulling this class back along a transverse section \(s \colon M \to E\) gives the Euler class, represented by the zero locus \(s^{-1}(0)\).

The higher characteristic classes can be described similarly by replacing the zero locus of one section by a linear-dependence locus.
For a rank \(r\) real or complex vector bundle \(E\), consider \(n+1\) sections of \(E\).
The locus where these sections become linearly dependent is generally singular, but it has a natural smooth resolution.
Over \(E^{\oplus(n+1)}\), define
\[
\chi_{r-n}(E) =
\Big\{ \big((v_0,\ldots,v_n),\, [a_0:\cdots:a_n]\big) \,\mid\, \textstyle \sum_i a_i v_i = 0 \Big\},
\]
where \([a_0:\cdots:a_n]\) lies in \(\mathbb{R}P^n\) in the real case and in \(\mathbb{C}P^n\) in the complex case.
The projection
\[
\chi_{r-n}(E) \to E^{\oplus(n+1)}
\]
has image the linear-dependence locus and defines a geometric cochain of codimension \(r-n\) in the real case, respectively complex codimension \(r-n\) in the complex case.
It represents the Stiefel--Whitney class \(w_{r-n}(E)\) with \(\mathbb{F}_2\)-coefficients in the real case and the Chern class \(c_{r-n}(E)\) with integral coefficients in the complex case.